\documentclass[a4paper,landscape,8pt]{extarticle}
\usepackage[landscape,margin=0.7cm,top=0.5cm,bottom=0.7cm]{geometry}
\usepackage[utf8]{inputenc}
\usepackage[T1]{fontenc}
\usepackage[ngerman]{babel}
\usepackage{helvet}
\renewcommand{\familydefault}{\sfdefault}
\usepackage{xcolor}
\usepackage{tikz}
\usetikzlibrary{positioning, shapes.geometric}
\usepackage{enumitem}
\usepackage{multicol}
\usepackage{array}
\usepackage{fancyhdr}
\usepackage{amssymb}

% Farben
\definecolor{homeblue}{HTML}{007AFF}
\definecolor{appgreen}{HTML}{34C759}
\definecolor{settingsgray}{HTML}{8E8E93}
\definecolor{gestureorange}{HTML}{FF9500}
\definecolor{lightgray}{HTML}{F5F5F5}
\definecolor{bordergray}{HTML}{BBBBBB}
\definecolor{warnred}{HTML}{C62828}

\pagestyle{fancy}
\fancyhf{}
\renewcommand{\headrulewidth}{0pt}
\fancyfoot[C]{\footnotesize Seite \thepage{} von 3}

\setlength{\parindent}{0pt}
\setlength{\parskip}{2pt}

% Kompakte Listen
\setlist{nosep, leftmargin=1em, itemsep=0pt, topsep=1pt}

\begin{document}

% ============================================================================
% SEITE 1: INFORMATIONSBLATT
% ============================================================================

\noindent{\Large\bfseries Dein iPhone kennenlernen: Bildschirm, Apps \& Gesten} \hfill
{\footnotesize\textbf{Lernziele:} Home-Bildschirm verstehen | Apps finden | Grundgesten beherrschen}

\vspace{1mm}
\hrule
\vspace{1mm}

% Drei-Spalten-Layout
\begin{minipage}[t]{0.32\linewidth}
\begin{center}
\colorbox{homeblue!10}{%
\begin{minipage}{0.95\linewidth}
\vspace{2mm}
\begin{center}
{\Large\bfseries\textcolor{homeblue}{Der Home-Bildschirm}}\\[1pt]
{\footnotesize\itshape Dein digitaler Schreibtisch}
\end{center}
\vspace{2mm}
\end{minipage}%
}
\end{center}

\vspace{1mm}

\textbf{Was ist das?} Der Bildschirm, den du siehst, wenn du dein iPhone entsperrst.

\textbf{Analogie:} Wie dein \textbf{Schreibtisch} -- hier liegen deine Werkzeuge (Apps) griffbereit.

\textbf{Die Elemente:}
\begin{itemize}
\item \textbf{Apps} = kleine Symbole (Programme)
\item \textbf{Dock} = die 4 Apps ganz unten (immer sichtbar)
\item \textbf{Punkte unten} = weitere Seiten (wischen!)
\item \textbf{Statusleiste oben} = Uhrzeit, WLAN, Akku
\end{itemize}

\textbf{Mehrere Seiten:}\\
Wische nach \textbf{links}, um weitere App-Seiten zu sehen. Wie \textbf{Schubladen} im Schreibtisch.

\textbf{App-Mediathek:}\\
Ganz nach links wischen: Hier sind \textbf{alle} Apps alphabetisch sortiert. Praktisch, wenn du eine App nicht findest!

\end{minipage}%
\hfill%
\begin{minipage}[t]{0.32\linewidth}
\begin{center}
\colorbox{appgreen!10}{%
\begin{minipage}{0.95\linewidth}
\vspace{2mm}
\begin{center}
{\Large\bfseries\textcolor{appgreen}{Apps verstehen}}\\[1pt]
{\footnotesize\itshape Die Werkzeuge auf deinem iPhone}
\end{center}
\vspace{2mm}
\end{minipage}%
}
\end{center}

\vspace{1mm}

\textbf{Was ist eine App?} Ein \textbf{Programm} für eine bestimmte Aufgabe.

\textbf{Analogie:} Wie \textbf{Werkzeuge} -- jedes für einen Zweck: Hammer, Schraubenzieher, Zange...

\textbf{Wichtige vorinstallierte Apps:}
\begin{itemize}
\item \textbf{Telefon} (grün) = Anrufen
\item \textbf{Nachrichten} (grün) = SMS/iMessage
\item \textbf{Safari} (Kompass) = Internet
\item \textbf{Fotos} (bunte Blume) = Bilder ansehen
\item \textbf{Kamera} = Fotografieren
\item \textbf{Einstellungen} (Zahnrad) = Alles einstellen
\end{itemize}

\textbf{App starten:} Einmal \textbf{tippen} auf das Symbol.

\textbf{App schließen:}
\begin{enumerate}
\item Von unten nach oben wischen und \textbf{halten}
\item App nach oben \textbf{wegwischen}
\end{enumerate}

\textbf{Neue Apps:} Kommen aus dem \textbf{App Store} (blaues A-Symbol).

\end{minipage}%
\hfill%
\begin{minipage}[t]{0.32\linewidth}
\begin{center}
\colorbox{gestureorange!10}{%
\begin{minipage}{0.95\linewidth}
\vspace{2mm}
\begin{center}
{\Large\bfseries\textcolor{gestureorange}{Die Grundgesten}}\\[1pt]
{\footnotesize\itshape So bedienst du dein iPhone}
\end{center}
\vspace{2mm}
\end{minipage}%
}
\end{center}

\vspace{1mm}

\textbf{Was sind Gesten?} \textbf{Fingerbewegungen} auf dem Bildschirm.

\textbf{Analogie:} Wie \textbf{Handgriffe} -- einmal gelernt, machst du sie automatisch.

\textbf{Die wichtigsten Gesten:}

\begin{tabular}{@{}p{2.5cm}p{5cm}@{}}
\textbf{Tippen} & Kurz berühren = Auswählen/Öffnen \\[2mm]
\textbf{Halten} & Länger drücken = Mehr Optionen \\[2mm]
\textbf{Wischen} & Finger ziehen = Seiten wechseln \\[2mm]
\textbf{Scrollen} & Nach oben/unten wischen = Lesen \\[2mm]
\textbf{Ziehen} & Festhalten + bewegen = Verschieben \\[2mm]
\textbf{Zoomen} & Zwei Finger auseinander = Vergrößern \\
\end{tabular}

\vspace{2mm}

\textbf{Zurück zum Home-Bildschirm:}\\
Von \textbf{ganz unten} nach \textbf{oben wischen}.\\
{\footnotesize (Bei älteren iPhones: Home-Taste drücken)}

\end{minipage}

\vspace{1mm}

% Zusammenfassung
\begin{center}
\fcolorbox{bordergray}{lightgray}{%
\begin{minipage}{0.98\linewidth}
\centering
\vspace{1.5mm}
\textbf{Merke:}
\textcolor{homeblue}{\textbf{Home-Bildschirm}} = dein Schreibtisch ~$|$~
\textcolor{appgreen}{\textbf{Apps}} = deine Werkzeuge ~$|$~
\textcolor{gestureorange}{\textbf{Gesten}} = deine Handgriffe ~$|$~
\textbf{Dock} = die 4 wichtigsten immer unten
\vspace{1.5mm}
\end{minipage}%
}
\end{center}

\newpage

% ============================================================================
% SEITE 2: EINSTELLUNGEN & KONTROLLZENTRUM
% ============================================================================

\begin{center}
{\LARGE\bfseries Die wichtigsten Einstellungen \& das Kontrollzentrum}\\[3pt]
{\normalsize Schnellzugriff auf häufig gebrauchte Funktionen}
\end{center}

\vspace{2mm}
\hrule
\vspace{3mm}

\begin{multicols}{2}

\section*{\textcolor{settingsgray}{Das Kontrollzentrum}}

\textbf{Was ist das?} Eine \textbf{Schnellauswahl} für häufige Einstellungen -- ohne die Einstellungen-App zu öffnen.

\textbf{Analogie:} Wie die \textbf{Schalttafel} im Auto -- Licht, Heizung, Radio direkt erreichbar.

\textbf{So öffnest du es:}\\
Von der \textbf{oberen rechten Ecke} nach \textbf{unten wischen}.

\vspace{2mm}

\textbf{Was findest du dort?}

\begin{tabular}{@{}cl@{}}
\tikz\draw[fill=homeblue] (0,0) rectangle (0.4,0.4); & \textbf{WLAN} -- Ein/Aus, Netzwerk wählen \\[2mm]
\tikz\draw[fill=homeblue] (0,0) rectangle (0.4,0.4); & \textbf{Bluetooth} -- Kopfhörer, Geräte verbinden \\[2mm]
\tikz\draw[fill=homeblue] (0,0) rectangle (0.4,0.4); & \textbf{Flugmodus} -- Alles aus (im Flugzeug) \\[2mm]
\tikz\draw[fill=gestureorange] (0,0) rectangle (0.4,0.4); & \textbf{Helligkeit} -- Bildschirm heller/dunkler \\[2mm]
\tikz\draw[fill=gestureorange] (0,0) rectangle (0.4,0.4); & \textbf{Lautstärke} -- Töne lauter/leiser \\[2mm]
\tikz\draw[fill=appgreen] (0,0) rectangle (0.4,0.4); & \textbf{Taschenlampe} -- Blitz als Licht nutzen \\[2mm]
\tikz\draw[fill=appgreen] (0,0) rectangle (0.4,0.4); & \textbf{Kamera} -- Schnellstart \\[2mm]
\tikz\draw[fill=warnred] (0,0) rectangle (0.4,0.4); & \textbf{Nicht stören} -- Anrufe stumm \\
\end{tabular}

\vspace{3mm}

\textbf{Kontrollzentrum schließen:}\\
Irgendwo außerhalb tippen oder nach oben wischen.

\vspace{4mm}

\section*{\textcolor{settingsgray}{Die Einstellungen-App}}

\textbf{Was ist das?} Die \textbf{Schaltzentrale} für \textit{alles} auf deinem iPhone.

\textbf{Symbol:} Graues Zahnrad \tikz\draw[fill=settingsgray] (0,0) circle (0.15);

\textbf{Analogie:} Wie der \textbf{Sicherungskasten} zu Hause -- hier stellst du alles ein.

\columnbreak

\section*{Die wichtigsten Einstellungen}

\begin{tabular}{@{}p{3.5cm}p{6cm}@{}}
\textbf{[Dein Name]} (oben) & Apple-ID, iCloud, Abos \\[2mm]
\textbf{WLAN} & Netzwerk auswählen/verbinden \\[2mm]
\textbf{Bluetooth} & Kopfhörer, Lautsprecher koppeln \\[2mm]
\textbf{Mitteilungen} & Welche App darf stören? \\[2mm]
\textbf{Töne \& Haptik} & Klingelton, Lautstärke \\[2mm]
\textbf{Anzeige \& Helligkeit} & Schriftgröße, Nachtmodus \\[2mm]
\textbf{Allgemein} & Software-Update, Speicher \\[2mm]
\textbf{Bedienungshilfen} & Größerer Text, Vorlesen \\[2mm]
\end{tabular}

\vspace{4mm}

\begin{center}
\fcolorbox{homeblue}{homeblue!8}{%
\begin{minipage}{0.9\linewidth}
\vspace{2mm}
\textbf{Tipp: Schrift vergrößern}\\[2pt]
Einstellungen → Anzeige \& Helligkeit → Textgröße\\
Oder: Einstellungen → Bedienungshilfen → Anzeige \& Textgröße → Größerer Text
\vspace{2mm}
\end{minipage}}
\end{center}

\vspace{4mm}

\begin{center}
\fcolorbox{warnred}{warnred!8}{%
\begin{minipage}{0.9\linewidth}
\vspace{2mm}
\textbf{Vorsicht in den Einstellungen!}\\[2pt]
Ändere nur, was du verstehst. Bei Unsicherheit: Frag jemanden!\\
\textbf{Zurücksetzen} (ganz unten in Allgemein) = löscht alles!
\vspace{2mm}
\end{minipage}}
\end{center}

\vspace{4mm}

\section*{Suchen auf dem iPhone}

Von der \textbf{Mitte des Home-Bildschirms nach unten wischen} = Suchfeld.

Hier kannst du suchen nach:
\begin{itemize}
\item Apps (z.B. \glqq Kamera\grqq{})
\item Kontakten (z.B. \glqq Mama\grqq{})
\item Einstellungen (z.B. \glqq WLAN\grqq{})
\item Nachrichten, Notizen, allem!
\end{itemize}

\end{multicols}

\newpage

% ============================================================================
% SEITE 3: ÜBUNGEN & CHECKLISTE
% ============================================================================

\begin{center}
{\LARGE\bfseries Übungen \& Checkliste}\\[3pt]
{\normalsize Teste dein Wissen und übe die Gesten}
\end{center}

\vspace{2mm}
\hrule
\vspace{3mm}

\begin{multicols}{2}

\textbf{\large Teil A: Begriffe zuordnen}

Was ist was? Verbinde mit Linien oder schreibe die Nummer.

\vspace{2mm}

\renewcommand{\arraystretch}{1.8}
\begin{tabular}{@{}p{5cm}p{1cm}p{5cm}@{}}
1. Home-Bildschirm & & A. Schnellauswahl für WLAN, Helligkeit \\
2. Dock & & B. Alle Apps alphabetisch sortiert \\
3. Kontrollzentrum & & C. Die 4 Apps ganz unten \\
4. App-Mediathek & & D. Dein digitaler Schreibtisch \\
5. Einstellungen & & E. Die Schaltzentrale (Zahnrad) \\
\end{tabular}
\renewcommand{\arraystretch}{1}

\vspace{5mm}

\textbf{\large Teil B: Gesten üben}

Führe diese Aktionen auf deinem iPhone aus und hake ab:

\vspace{2mm}

\begin{itemize}[label=$\square$, itemsep=4pt]
\item Öffne eine App durch \textbf{Tippen}
\item Gehe zurück zum Home-Bildschirm (von unten wischen)
\item Öffne das \textbf{Kontrollzentrum} (von oben rechts wischen)
\item Schalte die \textbf{Taschenlampe} ein und wieder aus
\item Schließe das Kontrollzentrum
\item Wische zur \textbf{nächsten Seite} des Home-Bildschirms
\item Öffne die \textbf{App-Mediathek} (ganz nach links wischen)
\item Öffne die \textbf{Suche} (von der Mitte nach unten wischen)
\item Suche nach \glqq Einstellungen\grqq{} und öffne die App
\end{itemize}

\vspace{5mm}

\textbf{\large Teil C: Was machst du?}

\vspace{2mm}

\textbf{Situation 1:} Du willst das WLAN schnell ein- oder ausschalten.

\rule{7cm}{0.4pt}

\vspace{3mm}

\textbf{Situation 2:} Du findest eine App nicht auf dem Home-Bildschirm.

\rule{7cm}{0.4pt}

\vspace{3mm}

\textbf{Situation 3:} Die Schrift auf dem iPhone ist zu klein.

\rule{7cm}{0.4pt}

\columnbreak

\textbf{\large Checkliste: Das kann ich jetzt}

\vspace{2mm}

\begin{center}
\fcolorbox{homeblue}{homeblue!8}{%
\begin{minipage}{0.95\linewidth}
\vspace{2mm}
\textbf{\textcolor{homeblue}{Home-Bildschirm}}
\begin{itemize}[label=$\square$, itemsep=3pt]
\item Ich verstehe, dass der Home-Bildschirm mein \textbf{Schreibtisch} ist
\item Ich weiß, was das \textbf{Dock} ist (4 Apps unten)
\item Ich kann zu anderen \textbf{Seiten wischen}
\item Ich kann die \textbf{App-Mediathek} öffnen
\item Ich kann die \textbf{Suche} öffnen
\end{itemize}
\vspace{2mm}
\end{minipage}}
\end{center}

\vspace{3mm}

\begin{center}
\fcolorbox{appgreen}{appgreen!8}{%
\begin{minipage}{0.95\linewidth}
\vspace{2mm}
\textbf{\textcolor{appgreen}{Apps \& Gesten}}
\begin{itemize}[label=$\square$, itemsep=3pt]
\item Ich kann eine App \textbf{öffnen} (tippen)
\item Ich kann zum Home-Bildschirm \textbf{zurück} (von unten wischen)
\item Ich kann in Apps \textbf{scrollen} (auf/ab wischen)
\item Ich kann Bilder \textbf{zoomen} (zwei Finger auseinander)
\end{itemize}
\vspace{2mm}
\end{minipage}}
\end{center}

\vspace{3mm}

\begin{center}
\fcolorbox{settingsgray}{lightgray}{%
\begin{minipage}{0.95\linewidth}
\vspace{2mm}
\textbf{\textcolor{settingsgray}{Einstellungen}}
\begin{itemize}[label=$\square$, itemsep=3pt]
\item Ich kann das \textbf{Kontrollzentrum} öffnen
\item Ich weiß, wo die \textbf{Einstellungen-App} ist
\item Ich kann die \textbf{Schriftgröße} ändern
\item Ich weiß: Bei \glqq Zurücksetzen\grqq{} = \textbf{Vorsicht!}
\end{itemize}
\vspace{2mm}
\end{minipage}}
\end{center}

\vspace{4mm}

\begin{center}
\fcolorbox{bordergray}{white}{%
\begin{minipage}{0.95\linewidth}
\vspace{2mm}
\textbf{Meine Notizen}\\[3pt]
Meine wichtigsten Apps im Dock:\\[2mm]
\dotfill\\[3mm]
\dotfill\\[3mm]
Was ich noch üben möchte:\\[2mm]
\dotfill
\vspace{2mm}
\end{minipage}}
\end{center}

\end{multicols}

\vfill

\begin{center}
\footnotesize\textit{Stand: \today{} -- Übung macht den Meister! Probiere die Gesten immer wieder aus.}
\end{center}

\end{document}
